\section{Source paths now generate the shared decoder proof}
\label{sec:decoder-source-paths}

The incoming checked-decoder packet already supplied a sound finite-window
executor, a theorem composing finite chains of restored windows, and a concrete
certificate for K215. Its concrete data comprise 281 distinct local blocks,
234 context roles, 371 internal role edges, and 97 primitive-entry edges.
Those data and their three Lean modules are reused unchanged here. The
remaining gap was semantic: a caller still had to supply every window premise
of a \texttt{Chain}, and the identification of the certificate's source with
the upstream source had only been checked outside the kernel.

The new component closes both gaps. It proves that the actual source path
provides all the window facts, and it identifies the certificate's expanded
source with the literal upstream source by kernel reduction. It thus promotes
the shared decoder construction from local block soundness to a theorem about
complete decoder traversal of the repository's source program.

\subsection{Exact theorem and hypotheses}

Write $\operatorname{Walk}_{215}(P)=(e,d)$ for the repository's
\texttt{RegisterPath.walk} on \texttt{KSource215.p215}, with the same
$\texttt{nnodes}+1$ fuel and root used upstream. Let
$Q=\operatorname{consume}(P,d)$, so $|Q|=d$; this operation incorporates both
implicit zero padding and truncation of an instruction-pointer input.
Let $t:\mathbb Z\to\{0,1\}$ be an arbitrary complete tape, and let $a$ be the
initial head position. Assume
\[
 t(a-1)=1,\qquad t(a+i)=Q_i\quad(0\leq i<d).
\]
No condition is imposed on the other tape cells. The theorem
\texttt{DecoderSourceBinding.}\allowbreak\texttt{upstream\_walk\_reaches} proves the existence of
an execution length $N$, an entry state $q$, a scanned bit $b$, and the matching
primitive operation and argument, such that
\[
 d-1\leq N\leq 7(d-1),\qquad
 \operatorname{run}_{K215}^{N}(84,a,t)
   =(q,a+d-1,t).
\]
In particular the entire tape is restored, including its unrestricted
exterior. The final bit is $b=t(a+d-1)$. A certified role binds this state and
bit to the exact source primitive $e$ and to the literal K215 action selected
there. The execution ends just before that primitive action is performed.

The result is uniform in $P$, $e$, $d$, $a$, and $t$. It is not limited to the
blank-tape startup or a finite collection of sampled runs. Its source-word
hypothesis is the natural premise at a decoder entry; the existing register
representation and cleanliness invariant are responsible for establishing
that premise when this result is connected to the whole interpreter.

\subsection{How the proof composes}

The proof introduces two elementary tape predicates.
\texttt{Reads $t$ $a$ $w$} says that the list $w$ occurs starting at position
$a$. \texttt{Before $t$ $a$ $s$} says that the context $s$ ends immediately
before $a$. The generic list lemmas show that appending the currently scanned
bit extends this context, taking a suffix preserves its right alignment, and
a right-aligned context implies the overlay equality required by a local
block. Thus a block's finite-window hypothesis is derived from the source
word itself.

A role stores a source node, a decoder state, and a context of at most two
preceding bits. The source-path induction has two cases. For an internal
edge, the structural certificate identifies the child role, its updated
context, and the block to execute. The just-proved tape lemmas establish that
block's window premise. The induction hypothesis supplies the child's chain,
and \texttt{Chain.more} prepends the certified block. For a primitive edge,
zero internal steps remain; the certificate provides the exact native
transition row and its source leaf. This proves
\texttt{source\_path\_chain}. Applying the incoming
\texttt{chain\_sound} theorem gives the full-tape execution.

There is no duplicate concrete structural check. The generic theorem
\texttt{edgeSpec\_of\_check} extracts the semantic facts from the existing
Boolean \texttt{edgeCheck}. The theorem \texttt{edges\_of\_structure}
then derives them for every role from the already checked
\texttt{K215Decoder.structure\_checked}. The inherited
\texttt{all\_blocks\_checked} theorem supplies all local execution facts.
A small additional finite duration check verifies that every block takes
between one and seven steps. Induction on a chain proves the displayed
bounds, since each internal block advances the restored head by exactly one
source bit.

\subsection{Kernel binding to the upstream source}

The inherited certificate represents its source as 202 nodes: 179 internal
binary nodes and 23 distinct primitive leaves. The upstream register source
uses a 179-node DAG whose edges contain the primitive operations directly.
Their numbering conventions therefore differ. The adapter defines an explicit
expansion of the certificate nodes into the upstream \texttt{RMBisim.Tree}
type, translating the five operation tags into the corresponding
\texttt{RegisterPath.Edge} constructors.

The exact identity
\[
 \operatorname{sourceTree}(\texttt{sourceNodes},203,\texttt{sourceRoot})
   =\texttt{KSource215.t215}
\]
is proved by \texttt{rfl}. It is a kernel-checked definitional equality, with
no assumed correspondence and no external hash used as a proof premise.
The generic theorem \texttt{source\_path\_of\_leaf} converts every
upstream \texttt{SiteCoverage.LeafAt} witness through this identity into a
source-DAG path. Finally the existing \texttt{twalk\_leafAt},
\texttt{walk\_root}, and consumed-prefix equality connect the theorem to
\texttt{RegisterPath.walk}. The pinned source has 4,751 primitive occurrences,
with path depths from 3 to 47, as also recorded by the supplied source audit.

\subsection{Verification and remaining integration}

The portable component contains two new Lean modules, three unchanged
incoming decoder modules, and ten unchanged upstream source modules. Its
reproduction driver compiles this 15-module closure from source with official
Lean 4.32.0; Mathlib and prebuilt repository objects are unnecessary. The
source index identifies each inherited file and binds its bytes to repository
commit \texttt{aa66d0a33f}. Compiler receipts and logs accompany the component. The delivered verification
combines successful builds in the same session, with exact source, object,
and imported-object hashes; the two new modules were recompiled against that
closure. A subsequent single fresh replay was interrupted by an operating-system
termination while rebuilding the already verified source-data module, so it is
recorded as a resource failure rather than a successful complete replay.
The literal source-tree identity has no axioms. Both the complete source-path
and upstream-walk endpoints report only \texttt{propext},
\texttt{Classical.choice}, and \texttt{Quot.sound}; no admitted theorem or
native evaluator is introduced.

This closes the source-path and source-expansion obligations of the incoming
shared-decoder proposal. To replace \texttt{KSites215}, the next step is to
compose this theorem with context-parametric primitive-operation and pointer
return contracts. The existing 97 entry bindings identify the required
interfaces; they do not by themselves prove those operations. An adapter
between the native integer-tape presentation here and the repository's
renamed K230-label zipper presentation must also be supplied. In particular,
\texttt{KSim.Covered} currently requires a successful per-leaf Boolean check;
a reusable semantic-coverage interface would allow the shared traversal and
primitive contracts to feed its simulation theorem directly.

The benefit is a precise change in proof organization: internal decoding is
proved through 234 reusable contexts rather than separately for 4,750 internal
source occurrences. The full endpoint's compilation-time improvement remains
to be measured after primitive and return integration. The machine itself is
unchanged, and this decoder result does not constitute a 214-state
construction.
